\section{Introduction}
\label{sec:introduction}

A paper reports that reinforcement learning post-training moved a benchmark by two points. Whether
that is a result depends on a number nobody has: how much the same procedure would have moved it
had the seed been different. The field is aware of the problem and has documented it
\citep{hochlehnert2025sober}, and the only instrument it has is replication -- run the recipe $N$
times and look at the spread. At the compute scale RLVR now operates at, that instrument is close to
unaffordable, and so the number usually goes unmeasured.

This paper is about what happens between two runs of the same recipe, and it starts from an
observation that turns out to be wrong in a useful way.

If gradient noise perturbs each update and perturbations accumulate, two runs that begin from the
same policy execute a random walk away from each other. Their divergence grows linearly in the
number of updates, a long run is less reproducible than a short one, and the only defence is a
bigger batch. That picture is implicit in how run-to-run variance is discussed, and it is the
natural reading of a noisy gradient.

It does not survive measurement. Two runs which differ only in the randomness of their rollouts
diverge for a few tens of updates and then stop. Across the settings we measure, the log-log slope
of divergence against update count is at most $+0.26$ against the $+1$ a random walk requires, and
usually negative. The runs are climbing the same objective, and its curvature pulls them back
together as fast as the noise pushes them apart. We call this \emph{seed reconvergence}, and its
fixed point is what an error bar on a single run should be built from.

\paragraph{Contributions.}
\begin{itemize}
  \item \textbf{Seed divergence is stationary, not accumulating} (Section~\ref{sec:seeds},
  Section~\ref{sec:experiments}). We measure it on transformers trained from scratch across a
  tenfold range of batch sizes and step sizes, and on a pretrained model, and in no setting that is
  still learning does it accumulate.

  \item \textbf{A law with two measurable factors.}
  $\KL_\infty = \tau_c D_{\mathrm{noise}}$ (Proposition~\ref{prop:stationary}): the diffusive drift
  injected per update, times how long a perturbation survives. It predicts a $P^{-1/2}D^{+1/2}$
  scaling in the batch and the drift, measured at $P^{-0.45}D^{0.35}$ with $R^2 = 0.94$, and an
  outcome spread going as $\sqrt{\KL_\infty}$.

  \item \textbf{An error bar from one run.} $D_{\mathrm{noise}}$ comes from gradients a trainer is
  already accumulating (Section~\ref{sec:estimation}); $\tau_c$ comes from forking the run for a few
  tens of updates rather than repeating it. The alternative currently in use costs $N$ full runs.

  \item \textbf{Where it fails, and why that matters.} A run held to a fixed drift target must raise
  its step size as the objective saturates, because $\E_x[p(1-p)] \to 0$ takes the gradient with it.
  Runs that cross that point leave the contracting regime and separate. The failure is a property of
  drift-targeted control rather than of the law, and it is a caution about a technique this paper
  otherwise recommends.

  \item \textbf{The machinery underneath}, which is of independent use: an exact finite-$G$ mean and
  covariance for the advantage estimators in current use under binary rewards
  (Propositions~\ref{prop:mean} and \ref{prop:cov}). Two things fall out. The estimators differ only
  through a weight $\lambda(p,G)$ on prompt difficulty, so those with a $p$-independent $\lambda$ are
  the same estimator up to a step size. And the gradient noise separates into a between-prompt term
  and a within-prompt term, which is what makes $D_{\mathrm{noise}}$ measurable at all, and which
  incidentally fixes the group size (Theorem~\ref{thm:split}).
\end{itemize}

Every experiment here runs on one laptop. The claim under test is that a small measurement replaces
a large one, and a study that needed a cluster to take the measurement would be a poor advertisement
for it.
